%% ACM article starter — acmart document class.
%%
%% The acmart class itself is NOT included in this download: it ships with
%% every modern TeX distribution (and with latextemplates.com's compiler), so
%% bundling a copy would only risk pinning you to an older version than your
%% own system already has.
%%
%% Class documentation: texdoc acmart   |   https://ctan.org/pkg/acmart
%% Licence: acmart is LPPL-1.3, (c) Association for Computing Machinery.
%%
%% Pick the format your venue asks for by changing the option below:
%%   acmsmall   ACM small trim journals (default for most TOG/TOCHI journals)
%%   acmlarge   ACM large trim journals
%%   sigconf    most ACM conference proceedings
%%   sigplan    SIGPLAN conferences
%%   manuscript submission/review copy, single column, line numbers
\documentclass[sigconf]{acmart}

%% acmart does NOT load cleveref itself — load it here if you want \Cref.
%% (Checked against the class rather than assumed.)
\usepackage{cleveref}

%% Conference/journal metadata. Your venue supplies these exact strings —
%% replace them; do not leave the placeholders in a submission.
\acmConference[Conference'XX]{ACM Conference}{June 03--05, 2026}{Woodstock, NY}
\acmBooktitle{ACM Conference (Conference'XX), June 03--05, 2026, Woodstock, NY}

\begin{document}

\title{Your Title Here}
\subtitle{An optional subtitle}

\author{Ada Lovelace}
\email{ada@example.edu}
\affiliation{%
  \institution{Institute of Analytical Engines}
  \city{London}
  \country{United Kingdom}
}

\author{Second Author}
\email{second@example.org}
\affiliation{%
  \institution{Another University}
  \city{Springfield}
  \country{USA}
}

\begin{abstract}
  Replace this with your abstract. A good abstract states the problem, what
  you did, and what you found --- in that order, and without citations.
\end{abstract}

\keywords{first keyword, second keyword, third keyword}

\maketitle

\section{Introduction}
Start here. Cite things like this~\cite{knuth1984texbook}, and cross-reference
sections like \Cref{sec:method}.

\section{Method}
\label{sec:method}
Describe what you did.

\begin{table}[t]
  \caption{A caption goes above the table.}
  \label{tab:example}
  \begin{tabular}{lrr}
    \toprule
    Condition & Count & Rate \\
    \midrule
    Baseline  & 120   & 0.42 \\
    Ours      & 118   & 0.71 \\
    \bottomrule
  \end{tabular}
\end{table}

\section{Results}
\Cref{tab:example} shows the numbers.

\section{Conclusion}
Say what it means.

\bibliographystyle{ACM-Reference-Format}
\bibliography{references}

\end{document}
